\chapter{Introduction}
\label{chap:introduction}

\placeholder{Rename this chapter and replace every example paragraph with your own work.}

\section{Background and context}
Start broadly enough for a reader outside your immediate specialisation. Define the
real-world or scholarly setting, introduce the central concepts, and support factual
claims with authoritative sources. For example, research projects often become harder
to maintain when their writing, evidence, figures and revision history are scattered
across unrelated files \citep{knuth1984texbook}.

\subsection{Moving from a broad setting to a precise focus}
A subsection develops one part of the chapter argument without creating a new chapter.
This fully synthetic example demonstrates the standard hierarchy automatically:
Chapter~1 contains Section~1.1, which contains Subsection~1.1.1. Students should use
headings only when they help a reader follow the logic of the real thesis.

\section{Problem statement and research gap}
State what is known, what remains unresolved, and why that unresolved issue matters.
Avoid presenting a broad topic as though it were automatically a research gap.

\CorrectionAdd{E1-01}{a}{Although prior work explains how structured authoring tools
support reproducible documents, there remains limited practical guidance connecting
large-thesis file organisation, transparent examiner revisions and multiple submission
outputs in one maintainable workflow. This thesis addresses that operational gap.}

\section{Aim, objectives and research questions}
The aim should describe the overall intellectual destination. Objectives are the
concrete steps used to reach it, while research questions define what the thesis will
answer. A student could adapt the following pattern:
\begin{quote}
\textbf{Aim:} To investigate how \placeholder{phenomenon} affects
\placeholder{population or system} in \placeholder{context}.
\end{quote}

\begin{enumerate}
  \item Review and synthesise the relevant literature.
  \item Design and apply an appropriate research method.
  \item Analyse the evidence and explain its implications.
\end{enumerate}

\begin{description}
  \item[RQ1] How does \placeholder{factor A} influence \placeholder{outcome B}?
  \item[RQ2] Under what conditions does that relationship change?
\end{description}

\section{Illustrative contributions}
\begin{table}[htbp]
  \centering
  \caption{Example contribution map}
  \label{tab:contribution-map}
  \begin{tabularx}{\textwidth}{@{}p{0.20\textwidth}XX@{}}
    \toprule
    \textbf{Contribution} & \textbf{Claim to be demonstrated} & \textbf{Evidence location} \\
    \midrule
    C1 & A clearly defined conceptual or empirical contribution. & \Cref{chap:literature,chap:results} \\
    C2 & A methodological or practical contribution. & \Cref{chap:methodology,chap:discussion} \\
    C3 & A set of implications grounded in the study evidence. & \Cref{chap:discussion,chap:conclusion} \\
    \bottomrule
  \end{tabularx}
\end{table}


Do not claim novelty merely because a technique is new to your project. Connect every
claimed contribution to evidence developed in later chapters.

\section{Thesis structure}
\Cref{chap:literature} synthesises prior research. \Cref{chap:methodology} explains the
research design, \Cref{chap:results} reports the evidence, and \Cref{chap:discussion}
interprets it. \Cref{chap:conclusion} answers the research questions and identifies
future work.
